{\small\setlength{\tabcolsep}{4pt}
\begin{longtable}{llrrrr}
\caption{Combined-target sensitivity without an IGA component}\label{tab:iga_definition}\\
\toprule Week & Initial group & Paired N & Recorded & No IGA & Changed \\ \midrule\endfirsthead
\toprule Week & Initial group & Paired N & Recorded & No IGA & Changed \\ \midrule\endhead
W4 & dupi & 39 & 18/39 & 14/39 & 4 \\
W4 & abro 100 mg & 34 & 26/34 & 26/34 & 0 \\
W4 & abro 200 mg & 17 & 13/17 & 13/17 & 0 \\
W12 & dupi & 48 & 20/48 & 18/48 & 2 \\
W12 & abro 100 mg & 34 & 21/34 & 17/34 & 4 \\
W12 & abro 200 mg & 30 & 19/30 & 19/30 & 0 \\
W24 & dupi & 36 & 26/36 & 26/36 & 0 \\
W24 & abro 100 mg & 24 & 22/24 & 22/24 & 0 \\
W24 & abro 200 mg & 23 & 20/23 & 20/23 & 0 \\
\bottomrule\end{longtable}
}
Both definitions are evaluated in the same patients within each row. Recorded denotes the reconstructed framework, not a verified consensus endpoint. No IGA removes only the IGA alternative from the physician domain, retaining EASI50 at W4, EASI75 at W12 and EASI75 or EASI$\leq5$ at W24; the patient domain is unchanged. This is a different, exploratory target definition, not a recoding of IGA.

{\small\begin{tabular}{llrr}\toprule Week & Initial group & Recorded, all evaluable & No IGA, all evaluable \\ \midrule
W4 & dupi & 18/39 & 14/39 \\
W4 & abro 100 mg & 26/34 & 26/34 \\
W4 & abro 200 mg & 13/17 & 13/18 \\
W12 & dupi & 20/48 & 18/48 \\
W12 & abro 100 mg & 21/34 & 17/34 \\
W12 & abro 200 mg & 19/30 & 19/30 \\
W24 & dupi & 26/36 & 26/36 \\
W24 & abro 100 mg & 22/24 & 22/24 \\
W24 & abro 200 mg & 20/23 & 20/23 \\
\bottomrule\end{tabular}}

The second block uses each definition\textquotesingle s own evaluable population; a change in denominator is not a change in the same patient\textquotesingle s response. Unknown components follow three-valued AND/OR rules and are never filled with clinical scores. Full transition counts, physician-domain results and exclusion of the 22 patients with recorded baseline IGA5 are supplied in the source tables. The protocol inclusion text specifies 0--4 but its IGA appendix defines grades 0--5, including grade 5 as very severe. All original grades are therefore retained; actual implementation requires confirmation.
